\documentclass[11pt]{article}
% The build script supplies \TechnicaTheme for the dark variant.
\providecommand{\TechnicaTheme}{light}
\edef\LoadTechnica{\noexpand\usepackage[theme=\TechnicaTheme]{technica}}
\LoadTechnica
\usepackage{hyperref}
\title{Technica specimen}
\author{Technical writing, in context}
\date{Light and dark theme reference}

\begin{document}
\maketitle

\section{A readable technical page}
A \term{translation unit} is the input to one compilation step. Open
\filepath{src/main.cpp}, run \command{clang++}, and inspect the
\api{std::vector} API. Press \key{Esc} to return to the editor.
These semantic commands keep the document's vocabulary separate from its
presentation. A footnote remains part of the reading rhythm.\footnote{This is
a specimen of footnote spacing and type size.}

\subsection{Structure and notation}
For readings $x_1,\ldots,x_n$, the arithmetic mean is
\[
    \bar{x} = \frac{1}{n}\sum_{i=1}^{n} x_i.
\]
\begin{enumerate}
\item Read the input and validate each value.
\item Compute the total, then divide by the number of readings.
\item Report the result with its units.
\end{enumerate}

\begin{table}[h]
\centering
\begin{tabular}{lr}
\hline
Reading & Value \\
\hline
First & 12 \\
Second & 18 \\
Third & 24 \\
\hline
\end{tabular}
\caption{A compact table with a muted caption.}
\end{table}

\begin{note}
The same source can be rendered with either theme. Links such as
\href{https://www.latex-project.org/}{the LaTeX Project} use the selected palette.
\end{note}

\newpage
\section{Code and terminal sessions}
Inline operators remain literal: \inlinecpp{ready || retry} and
\inlinesh{cat input | sort}. Percent signs survive: \inlinecode{100% complete}.

Delimiters allow braces: \inlinecpp!if (ready) { run(); }!.

They also allow paths: \inlinecode!C:\tmp\input.txt!.

\technicainputcode[listingopts={style=technica-cpp},label=code:readings]{sample.cpp}{examples/sample.cpp}

\begin{terminal}[title={Build and run}]
$ clang++ -std=c++20 sample.cpp -o sample
$ ./sample
Total: 54
\end{terminal}

\begin{codeblock}[listingopts={numbers=none}]{Wrapping a long line}
message = "This deliberately long configuration value checks that code wraps inside the available text width while preserving all of the words in the original input."
\end{codeblock}

\begin{advice}
Include code from its source file so the example stays aligned with the program.
Block \ref{code:readings} above is read directly from \filepath{examples/sample.cpp}.
\end{advice}

\newpage
\section{Callouts across page boundaries}
\begin{warning}
An empty input has no arithmetic mean. Check the count before dividing by it.
\end{warning}

\begin{explanation}
This longer explanation exercises page breaking. A callout should remain
readable when its contents exceed the space left on a page.

\newcount\specimenparagraph
\specimenparagraph=0
\loop
  \advance\specimenparagraph by 1
  \textbf{Step \the\specimenparagraph.}
  Readings belong to a sequence with an explicit order. Keep that order when
  presenting intermediate results so readers can match the explanation to the
  example. Name each quantity before using it, and carry its units through the
  calculation. This paragraph supplies enough prose to exercise line fitting,
  paragraph spacing, and the continuation of a callout on the next page.\par
\ifnum\specimenparagraph<9
\repeat
\end{explanation}

\subsection{After the continuation}
Normal prose resumes after the box. The final paragraph checks that body text,
spacing, and the theme's ink color are restored.
\end{document}
